\subsection{Approximate Computing come leva per la riduzione dei consumi}

Una direzione promettente per la riduzione del consumo energetico
in ambito FaaS è rappresentata dal paradigma dell'\emph{Approximate
Computing}, che consente di ridurre il costo computazionale
accettando, in modo controllato, una diminuzione dell'accuratezza
del risultato prodotto. Lavori come JouleGuard~\cite{jouleguard}
hanno mostrato come sia possibile progettare meccanismi di controllo
runtime capaci di rispettare un budget energetico prefissato,
regolando dinamicamente il livello di approssimazione dell'esecuzione.

Applicata al paradigma FaaS, questa idea si traduce nella possibilità
di associare a una stessa funzione logica più varianti implementative,
caratterizzate da diversi profili energetici e da diversi livelli
di accuratezza del risultato. La selezione della variante da eseguire
può quindi diventare una decisione dello scheduler, guidata da
informazioni energetiche aggiornate a runtime e dai vincoli
espressi dall'utente al momento dell'invocazione.
